\scanpage{151}
% archive sheet number: 150
% L1–5
$\rho_0$, qui fixe $\lambda s_0$ (puisqu'il fixe $s_0$ et commute à $\lambda$) et induit un autom direct de l'étoile de $\lambda s_0$ \deleted{\unsure{pour}} pour $P_1$, est un automorphisme direct pour $P_1$ i.e. on a (54), ce qui achève de prouver
% L6–10
\[
\begin{aligned}
&\text{\deleted{\unsure{les}} les caractères d'orientation \inserted{sur $G_0$} pour les structures}\\
&\text{de polyèdres }P_i\quad(i\in\Pi)\text{ sont les $=$s, i.e.}\\
&\text{il existe un (unique) sous-groupe }G_0^+\text{ de}\\
&G_0,\text{ tel que l'on ait}\\
&\qquad G_0^+=\mathrm{Aut}^+(\unsure{P_i|\Phi_i})\qquad\forall i\in\Pi\ .
\end{aligned}
\tag*{(56)}
\]
% L11
Cela achève en même temps de prouver (48.).

% L12–16
Considérons maintenant l'action (par transport de structure) de $G=\mathrm{Aut}(S,F)$ sur $\unsure{\vec\Pi|\vec\pi}$, il est immédiat que cette action \deleted{\unsure{commute}}\inserted{respecte} la \deleted{\unsure{celle de}} structure de $\Lambda_n^*$-torseur sur $\vec\Pi$, d'où on déduit un unique homomorphisme
% L17
\[
G\xrightarrow{\chi}(\mathbb Z/n\mathbb Z)^*\qquad
\text{le ``\uline{caractère}}\quad\substack{\strut\\\text{\uline{\unsure{cyclotomique|cyclotromique}}''}}
\tag*{(57)}
\]
% L18–19
caractérisé par la propriété
\[
\forall\,\underbrace{i\in\Pi,\ \omega\in\mu(i)}_{\text{d'où }\alpha=(i,\omega)\in\vec\Pi},
\ \inserted{g\in G,}\ \text{on a}\quad
\underbrace{g\overbrace{(i,\unsure{\omega|\lambda})}^{\alpha\in\vec\Pi}}_{\substack{\text{transport}\\\text{de structure}}}
=\underbrace{\chi(g).\overbrace{(i,\omega)}^\alpha}_{\substack{\text{opération}\\\text{arithmétique}\\\text{d'exposant}\\
\chi(g)\text{ sur la}\\\text{structure}\\P_i=(S,A_i,F_i)}}.
\tag*{(58)}
\]

% L20–22
Le caractère $\chi$ factorise le caractère $\chi':G\longrightarrow(\mathbb Z/n\mathbb Z)^*/\pm1\deleted{\unsure{\simeq Z_n}}=Z'_n\simeq Z$ (``\uline{caractère demi-cyclotomique}''), i.e. on a commutativité \unsure{du|de|dans}
